\subsection{Risk Minimization}
% Loss functions, empirical risk, expected risk, consistency, bias-variance tradeoff.

\begin{remark} \hlt{Optimisation}\\
A cost function $J(\bm{\theta})$ is reduced to improve some performance measure $P$.\\
Optimisation algorithms include some specialisation on specific structure of machine learning objective functions.\\
Cost function is average over training set,
\begin{equation*}
J(\bm{\theta}) = \mathbb{E}_{(\bm{x}, y) \sim \hat{p}_{\text{data}}} L(f(\bm{x}; \bm{\theta}), y) \nonumber
\end{equation*}
where $L$ is per-example loss function, $f(\bm{x}; \bm{\theta})$ is predicted output given input $\bm{x}$, $\hat{p}_{\text{data}}$ is empirical distribution. For supervised learning, $y$ is target output.\\
The goal is to usually to minimise objective function where expectation is taken on data generating distribution $p_{\text{data}}$ rather than over the finite training set:
\begin{equation*}
J^*(\bm{\theta}) = \mathbb{E}_{(\bm{x}, y) \sim p_{\text{data}}} L(f(\bm{x}; \bm{\theta}), y) \nonumber
\end{equation*}
Note that machine learning algorithm usually minimises a surrogate loss function, and halts when convergence criterion based on early  stopping is satisfied.
\end{remark}

\begin{definition} \hlt{Empirical Risk Minimisation}\\
To minimise expected loss on training set, where $\hat{p}(\bm{x}, y)$ is the empirical distribution. The empirical risk is
\begin{equation*}
\mathbb{E}_{\bm{x}, y \sim \hat{p}_{\text{data}}(\bm{x}, y)}[L(f(\bm{x}, \bm{\theta}), y)] = \frac{1}{m} \sum\limits_{i=1}^m L(f(\bm{x}^{(i)}; \bm{\theta}), y^{(i)}) \nonumber
\end{equation*}
where $m$ is number of training samples. Note that empirical risk minimisation is prone to overfitting, as models with high capacity may memorise the training set.
\end{definition}
